\begin{abstract}
Fine-grained action recognition requires distinguishing actions that look
nearly identical in single frames and differ mainly in how the body moves.
Two difficulties follow from this setting. Since all classes share a nearly
identical background, the action region must be located rather than inferred
from scene cues; and since the discriminating signal is motion, it must be
captured at two time scales, the short-term displacement of body parts over a
few frames and the long-term evolution of the whole sequence. We propose
\dgtdm, a three-branch framework that assigns these two requirements to
separate branches. A frozen CLIP encoder serves as a spatial expert, and a
spatio-temporal guidance (STG) branch injects its stable priors into a
lightweight 2D temporal encoder through bilateral window cross-attention and
adaptive gating, so that the encoder concentrates on the action region. The
core of the framework is a dual-granularity temporal dynamics module that
adds temporal modeling to the 2D backbone at the two time scales: the MoDeA
branch predicts dense cross-frame sampling offsets to follow short-term
body-part displacement, and the TSA branch maintains a long-term state
memory to model the sequence-level evolution. We evaluate the framework on
Diving48 and Something-Something V2 and use ablations to isolate the
contribution of each component.
\end{abstract}
